\documentclass[a4paper,10pt]{report}\input{../0.Préambule/Préambule_cours_prof}\begin{document}

\pagestyle{empty}
\newpage
\noindent NOM, Prénom et classe : \dotfill

\section*{Proportionnalité\hspace{3cm} Sujet A}
\addcontentsline{toc}{section}{Proportionnalité\hspace{3cm} Sujet A}

\noindent \textit{L'usage de la calculatrice est \textbf{autorisée}.} \hfill \textit{Durée : 55 minutes}

\noindent \textit{Le sujet est a rendre avec la copie.}

\paragraph{Exercice 1}: Question de cours\hfill \textbf{2 points}

\medskip
Lorsque l'on représente une situation de proportionnalité dans un repère, qu'obtient-on ?

\paragraph{Exercice 2}: \hfill \textbf{3 points}

\medskip
Sans justification, complète le tableau de proportionnalité ci-dessous.(On arrondira éventuellement les résultats au \textbf{centième})

\begin{center}
\renewcommand{\arraystretch}{2}
\begin{tabularx}{\linewidth}{|>{\columncolor{lightgray!20}}p{5cm}|*{4}{>{\centering \arraybackslash}X|}}
\hline
Taille d'un fichier (en Mo) &  & $3,75$ & $740$ &  \\\hline
Durée de téléchargement (en s) & $208$ & $44$ &  & $10$ \\\hline
\end{tabularx}
\end{center}

\paragraph{Exercice 3}: \hfill \textbf{6 points}

\medskip
\noindent \textit{La rédaction sera évaluée au même titre que les résultats, on prendra donc soin de justifier chacun de nos calculs.} 

\begin{enumerate}
\item J'ai acheté $12$m de ruban pour $6,40$\euro. Combien coûtent $7$m de ruban ?
\item J'ai utilisé $50$kg de semences pour un terrain de $\np{1 900}$m$^2$.
Quelle surface aurais-je pu ensemencer avec $90$kg de semences ?
\item En roulant à une vitesse moyenne de $72$km/h, quelle est la distance parcourue en $32$min ?
\end{enumerate}

\paragraph{Exercice 4}: \hfill \textbf{6 points}

\noindent \textit{La rédaction sera évaluée au même titre que les résultats, on prendra donc soin de justifier chacun de nos calculs.}

\medskip
Ce tableau indique la taille de Néo en fonction de son âge.
\begin{center}
\renewcommand{\arraystretch}{2}
\begin{tabularx}{\linewidth}{|>{\columncolor{lightgray!20}}p{4cm}|*{4}{>{\centering \arraybackslash}X|}}
\hline
Âge (en années) & $2$ & $5$ & $10$ & $12$ \\\hline
Taille (en cm) & $80$ & $110$ & $135$ & $160$ \\\hline
\end{tabularx}
\end{center}

\begin{enumerate}
\item Est-ce une situation de proportionnalité ?
\item En suivant les indication ci-dessous, représente graphiquement l'évolution de la taille de Néo en fonction de son âge. 

On prendra $1$ carreau pour $1$ an sur l'axe des abscisses.

On prendre $1$ carreau pour $10$ cm sur l'axe des ordonnées.
\item Peux-tu répondre à la question \textbf{{1.}} sans faire de calculs ? Justifie
\end{enumerate}

\paragraph{Exercice 5}: \hfill \textbf{3 points}

\medskip
\noindent \textit{Cet exercice est plus compliqué que les autres et ne rapporte pas beaucoup de points, je vous conseille de le traiter en dernier une fois tous les autres terminés.} 

\noindent \textit{La rédaction sera évaluée au même titre que les résultats, on prendra donc soin de justifier chacun de nos calculs.} 

\medskip
$20$ bûcherons ont abattu $156$ arbres en $3$ jours.

\medskip 
\noindent En travaillant au même rythme, combien d'arbres $4$ de ces bûcherons abattraient-ils en $15$ jours ?

\newpage
\noindent NOM, Prénom et classe : \dotfill

\section*{Proportionnalité\hspace{3cm} Sujet B}
\addcontentsline{toc}{section}{Proportionnalité\hspace{3cm} Sujet B}

\noindent \textit{L'usage de la calculatrice est \textbf{autorisée}.} \hfill \textit{Durée : 55 minutes}

\paragraph{Exercice 1}: Question de cours\hfill \textbf{2 points}

\medskip
Lorsque l'on représente une situation de proportionnalité dans un repère, qu'obtient-on ?

\paragraph{Exercice 2}: \hfill \textbf{3 points}

\medskip
Sans justification, complète le tableau de proportionnalité ci-dessous.(On arrondira éventuellement les résultats au \textbf{centième})

\begin{center}
\renewcommand{\arraystretch}{2}
\begin{tabularx}{\linewidth}{|>{\columncolor{lightgray!20}}p{5cm}|*{4}{>{\centering \arraybackslash}X|}}
\hline
Taille d'un fichier (en Mo) &  & $2,75$ & $740$ &  \\\hline
Durée de téléchargement (en s) & $208$ & $54$ &  & $10$ \\\hline
\end{tabularx}
\end{center}

\paragraph{Exercice 3}: \hfill \textbf{6 points}

\medskip
\noindent \textit{La rédaction sera évaluée au même titre que les résultats, on prendra donc soin de justifier chacun de nos calculs.} 

\begin{enumerate}
\item J'ai acheté $10$m de ruban pour $5,40$\euro. Combien coûtent $7$m de ruban ?
\item J'ai utilisé $60$kg de semences pour un terrain de $\np{1 600}$m$^2$.
Quelle surface aurais-je pu ensemencer avec $90$kg de semences ?
\item En roulant à une vitesse moyenne de $76$km/h, quelle est la distance parcourue en $25$min ?
\end{enumerate}

\paragraph{Exercice 4}: \hfill \textbf{6 points}

\noindent \textit{La rédaction sera évaluée au même titre que les résultats, on prendra donc soin de justifier chacun de nos calculs.}

\medskip
Ce tableau indique la taille de Noa en fonction de son âge.
\begin{center}
\renewcommand{\arraystretch}{2}
\begin{tabularx}{\linewidth}{|>{\columncolor{lightgray!20}}p{4cm}|*{4}{>{\centering \arraybackslash}X|}}
\hline
Âge (en années) & $2$ & $6$ & $11$ & $13$ \\\hline
Taille (en cm) & $80$ & $100$ & $125$ & $150$ \\\hline
\end{tabularx}
\end{center}

\begin{enumerate}
\item Est-ce une situation de proportionnalité ?
\item En suivant les indication ci-dessous, représente graphiquement l'évolution de la taille de Noa en fonction de son âge. 

On prendra $1$ carreau pour $1$ an sur l'axe des abscisses.

On prendre $1$ carreau pour $10$ cm sur l'axe des ordonnées.
\item Peux-tu répondre à la question \textbf{{1.}} sans faire de calculs ? Justifie
\end{enumerate}

\paragraph{Exercice 5}: \hfill \textbf{3 points}

\medskip
\noindent \textit{Cet exercice est plus compliqué que les autres et ne rapporte pas beaucoup de points, je vous conseille de le traiter en dernier une fois tous les autres terminés.} 

\noindent \textit{La rédaction sera évaluée au même titre que les résultats, on prendra donc soin de justifier chacun de nos calculs.} 

\medskip
$20$ bûcherons ont abattu $156$ arbres en $3$ jours.

\medskip 
\noindent En travaillant au même rythme, combien d'arbres $5$ de ces bûcherons abattraient-ils en $15$ jours ?
\end{document}
